\section{结论}
\label{sec:concl}

本文提出了一套多Agent自动化文献综述系统，解决了三个关键挑战：（1）手动文献筛选的可扩展性问题；（2）自动化分析的深度问题；（3）LLM生成内容的可靠性问题。

系统的核心创新包括：层级RQ驱动的流水线确保各阶段数据流的连贯性；三级筛选机制在吞吐量和精度之间取得平衡；GPU加速的语义聚类支持自适应参数调节；聚焦原理的报告生成策略通过多层反幻觉保障确保输出质量。

本系统的一个显著特点是强调公式级分析。与现有自动化综述工具产生浅层摘要不同，本系统通过精心设计的提示词强制要求每个方法必须附带核心公式及逐符号解释、从基本定律出发的推导链条、以及从第一性原理推导的理论性能极限。这种分析深度得益于层级RQ结构中包含的专门针对数学公式、物理模型和假设有效性的微观级问题。

系统集成五个免费学术数据库无需API密钥，支持命令行和Web界面两种交互方式，并通过检查点机制实现工作流中断恢复。GPU加速PCA、FAISS向量索引和并行LLM调用等性能优化使系统能在数分钟内处理数百篇论文。

未来工作将聚焦于：（1）引入全文PDF分析以克服当前仅依赖摘要的局限性，实现更深入的公式和推导提取；（2）扩展RQ树以支持不同研究领域的领域特定模板；（3）实现迭代优化机制，系统自动识别覆盖缺口并扩展检索范围；（4）开发自动化综述质量评估指标，包括公式准确性验证。
